\chapter{Формула Стокса.}

\begin{defn}
Пусть $G$ "--- ограниченная область в $\bbR^2_{uv}$. Отображение $F$ из $\overline{G}$ в $\bbR^3_{xyz}$, заданное векторным равенством $\vec{r} = \vec{r}(u, v)$, $(u, v) \in \overline{G}$, называется \textit{гладким\rindex{отображение!гладкое}} в $\overline{G}$, если оно непрерывно дифференцируемо в $\overline{G}$, причём $[\vec{r}_u, \vec{r}_v] \neq \vec{0}$ в $\overline{G}$ (т.е. векторы $\vec{r}_u$ и $\vec{r}_v$ не коллинеарны в $\overline{G}$).
\end{defn}

\begin{defn}
Пусть отображение $F$ биективно отображает замыкание $\overline{G}$ ограниченной области $G \subset \bbR^2_{uv}$ на множество $S \subset \bbR^3_{xyz}$, причём это отображение является гладким в $\overline{G}$. Тогда множество точек $S$ называется \textit{простой гладкой поверхностью\rindex{поверхность!простая гладкая}} (ПГП) в $\bbR^3$.
\end{defn}

\begin{exmpl}
График функции $z = f(x, y)$, непрерывно дифференцируемой в $\overline{G}$ ($G$ "--- ограниченная область в $\bbR^2_{xy}$), является ПГП.
\end{exmpl}

\begin{thm}[формула Стокса]
Пусть
\begin{enumerate}
\item[$1^\circ$.] вектор-функция $\vec{a} = (P, Q, R)^T$ непрерывно дифференцируема в области $G \subset \bbR^3_{xyz}$;
\item[$2^\circ$.] $S$ "--- дважды непрерывно дифференцируемая ПГП, заданная уравнением $\vec{r} = \vec{r}(u, v)$, $(u, v) \in \overline{D}$, где к области $D$ применима формула Грина, а граница $D$ "--- простая замкнутая кусочно-гладкая кривая $\gamma$, причём $S \subset G$;
\item[$3^\circ$.] ориентация $S$ согласована с ориентацией её края $\Gamma$.
\end{enumerate}
Тогда
$$
\int\limits_\Gamma (\vec{a}, d\vec{r}) = \iint\limits_S (\rot \vec{a}, d\vec{S}).
$$
\end{thm}

\begin{notion}
Под дважды непрерывной дифференцируемостью $S$ понимается дважды непрерывная дифференцируемость вектор-функции $\vec{r}(u, v)$ в $\overline{D}$.
\end{notion}

\begin{proof}
Пусть гладкая поверхность $S$ задана уравнением~\eqref{yaa14e1} и ориентирована своими параметрами, а граница $\partial D$ области $D$ состоит из одного или нескольких кусочно-гладких кусков вида $\gamma = \{u(t),v(t),t\in[a;b]\}$, каждый из которых параметром $t$ ориентирован положительно относительно области $D$. Тогда, как известно, граница $\partial S$ поверхности $S$ состоит из конечного числа кусочно-гладких кусков вида %недопечатнный параграф, но после как известно раньше было вот это --- (см.выше???)
$$
\Gamma = \{\vv{r}(u(t),v(t)),\quad t\in[a;b]\},
$$
ориентация каждого из которых параметром $t$ согласована с ориентацией поверхности $S$.

Для каждого куска $\Gamma\subset\partial S$ по формуле замены переменных
$$
\int\limits_{\Gamma} P\,dx = \int\limits_\gamma (Px'_u)\,du+(Px'_v)\,dv.
$$
Просуммировав по всем $\Gamma\subset\partial S$, получим
\begin{equation} \label{yaa14e5}
\int\limits_{\partial S} P\,dx = \int\limits_{\partial D} (Px'_u)\,du+(Px'_v)\,dv.
\end{equation}

При дополнительном условии, что функция $x(u,v)$ имеет непрерывную смешанную производную, интеграл по $\partial D$, согласно формуле Грина, равен следующему двойному интегралу:
$$
\iint\limits_D\left(\pd{(Px'_v)}{u} -\pd{(Px'_u)}{v}\right)\,du\,dv.
$$
Преобразуем подынтегральное выражение этого интеграла:
\begin{multline*}
\pd{(Px'_v)}{u} - \pd{(Px'_u)}{v}
=\left(\pd{P}{x}x'_u+\pd{P}{y}y'_u+\pd{P}{z}z'_u\right)x'_v-\\-\left(\pd{P}{x}x'_v+\pd{P}{y}y'_v+\pd{P}{z}z'_v\right)x'_u = \pd{P}{y}(y'_ux'_v-y'_vx'_u)+\pd{P}{z}(z'_ux'_v-z'_vx'_u)=\\=-\pd{P}{y}\cdot \pd{(x,y)}{(u,v)}+\pd{P}{z}\cdot\pd{(z,x)}{(u,v)}.
\end{multline*}
В результате получим равенство
\begin{equation} \label{yaa14e6}
\iint\limits_{\partial D} (Px'_u)\,du+(Px'_v)\,dv = \int\limits_D\left(\pd{P}{z}\cdot\pd{(z,x)}{(u,v)}-\pd{P}{y}\cdot \pd{(x,y)}{(u,v)}\right)\,du\,dv.
\end{equation}
Отметим, что оно справедливо и без предположения, что функция $x(u,v)$ имеет непрерывную смешанную производную. (Это утверждение можно доказать предельным переходом.)

Из равенств \eqref{yaa14e5}, \eqref{yaa14e6} и формулы для вычисления поверхностных интегралов следует, что
$$
\int\limits_{\partial S} P\,dx = \iint\limits_D \pd{P}{z}\cdot\pd{(z,x)}{(u,v)}du\,dv -\iint\limits_D \pd{P}{y}\cdot\pd{(x,y)}{(u,v)}du\,dv = \iint\limits_S \pd{P}{z}dz\,dx-\pd{P}{y}dx\,dy.
$$
Формула \eqref{yaa14e2} доказана. Формулы \eqref{yaa14e3}, \eqref{yaa14e4} доказываются аналогично.
\end{proof}

Имеет место следующее обобщение этой теоремы.

\begin{thmn}[формула Стокса, общий случай]
Пусть
\begin{enumerate}
\item[$1^\circ$.] вектор-функция $\vec{a} = (P, Q, R)^T$ непрерывно дифференцируема в области $G \subset \bbR^3_{xyz}$;
\item[$2^\circ$.] $S$ "--- ориентируемая КГП (см. определения 20.11 и 20.14) такая, что все ПГП $S_i$ дважды непрерывно дифференцируемы и задаются уравнениями $\vec{r} = \vec{r}(u, v)$, $(u, v) \in \overline{D}_i$, $i = 1, 2, \ldots, N$, где к областям $D_i$ применима формула Грина, а границы $D_i$ "--- простые замкнутые кусочно-гладкие кривые $\gamma_i$, причём $S \subset G$;
\item[$3^\circ$.] ориентация $S$ согласована с ориентацией её края $\Gamma$.
\end{enumerate}
Тогда
$$
\int\limits_\Gamma (\vec{a}, d\vec{r}) = \iint\limits_S (\rot \vec{a}, d\vec{S}).
$$
\end{thmn}

\begin{notion}
    Простая гладкая поверхность $S$ задаётся биективным отображением $F$ замыкания ограниченной плоской области $G \subset \bbR^2_{uv}$ в $\bbR^3_{xyz}$. Поэтому если граница $G$ является простой замкнутой кусочно-гладкой кривой в $\bbR^2_{uv}$, то край поверхности $S$ как образ $\partial G$ на поверхности является простой замкнутой кусочно-гладкой кривой в $\bbR^3_{xyz}$. Это следует из биективности отображения и из того, что образ гладкой (а значит, и кусочно-гладкой) кривой в $\overline{G}$ является гладкой (соответственно кусочно-гладкой) кривой на поверхности.
\end{notion}


